% Drawing for measurement `jfet-distortion`
% Stands for the deck lines:
%   VDD vdd 0 18 / RD vdd d {rd} / RS s 0 {rs}
%   RIN in g 10000 / VIN in 0 SIN(0 {amp} 1000)
%   J1 d g s DUT
%   .tran ... ten cycles
\documentclass[border=8pt]{standalone}
\usepackage[siunitx,RPvoltages,european]{circuitikz}
\ctikzset{bipoles/length=0.9cm, label/align=straight, inductors/scale=0.8}
\ctikzset{inductor=american}
\usepackage{amsmath}
\begin{document}
\begin{circuitikz}[font=\footnotesize, line width=0.5pt]
\draw (0,0) node[njfet](A){};
\node[anchor=west] at ($(A.D)!0.5!(A.S) + (0.55,0)$) {DUT};
\draw (A.S) to[short] ++(0,-0.3) coordinate(As) to[R, l_={$R_S$ ($I_D=I_{DSS}/3$)}] ++(0,-1.6) node[ground]{};
\draw (A.D) to[short] ++(0,0.3) coordinate(Ad) to[R, l={$R_D$ ($V_{DS}=\SI{9}{V}$)}] ++(0,1.5) node[vcc]{\SI{18}{V}};
\draw (A.G) to[short] ++(-0.4,0) to[R, l_={\SI{10}{k\ohm}}] ++(-1.6,0) to[short] ++(-0.4,0) to[vsourcesin, l_={$v_{in}$}, invert] ++(0,-1.9) node[ground]{};
\draw (Ad) to[short, -o] ++(1.2,0) node[right]{$v_{out}$};
\node[anchor=west, align=left] at (-4.6,3.3) {common source, unbypassed, 1 and \SI{3}{V} peak out};
\draw (9,0) node[njfet](B){};
\node[anchor=west] at ($(B.D)!0.5!(B.S) + (0.55,0)$) {DUT};
\draw (B.S) to[short] ++(0,-0.3) coordinate(Bs) to[R, l_={$R_S$ ($I_D=I_{DSS}/3$)}] ++(0,-1.6) node[ground]{};
\draw (B.D) to[short] ++(0,0.5) node[vcc]{\SI{18}{V}};
\draw (B.G) to[short] ++(-0.4,0) to[R, l_={\SI{10}{k\ohm}}] ++(-1.6,0) to[short] ++(-0.4,0) to[vsourcesin, l_={$v_{in}$}, invert] ++(0,-1.9) node[ground]{};
\draw (Bs) to[short] ++(0.7,0) to[C, l={\SI{100}{\micro F}}] ++(1.4,0) coordinate(Bl) to[R, l={\SI{10}{k\ohm}}] ++(0,-1.6) node[ground]{};
\draw (Bl) to[short, -o] ++(0.9,0) node[right]{$v_{out}$};
\node[anchor=west, align=left] at (6.5,3.3) {source follower, \SI{1}{V} peak out};
\end{circuitikz}
\end{document}
